% Generated draft: requires mathematical and source review.
\documentclass[11pt]{article}
\usepackage[margin=25mm]{geometry}
\usepackage{fontspec}
\setmainfont{Latin Modern Roman}
\usepackage{amsmath,amssymb,mathtools,unicode-math}
\setmathfont{Latin Modern Math}
\usepackage{enumitem,needspace,xurl,hyperref}
\setlength{\parindent}{0pt}
\setlength{\parskip}{5pt}
\setlength{\emergencystretch}{3em}
\newcommand{\N}{\mathbb{N}}
\newcommand{\Z}{\mathbb{Z}}
\newcommand{\Q}{\mathbb{Q}}
\newcommand{\R}{\mathbb{R}}
\newcommand{\C}{\mathbb{C}}
\newcommand{\F}{\mathbb{F}}
\newcommand{\A}{\mathbb{A}}
\newcommand{\PP}{\mathbb{P}}
\newcommand{\E}{\mathbb{E}}
\newcommand{\eps}{\varepsilon}
\newcommand{\dd}{\,\mathrm{d}}
\begin{document}
\section*{UnsolvedMath ID: 6800013}
\textbf{Source problem code:} catalogue entry
\subsection{Definitions and mathematical statement}
A coarse embedding \(f:(X,d_X)\to(Y,d_Y)\) is a map for which there are nondecreasing control functions \(\rho_-,\rho_+:[0,\infty)\to[0,\infty)\), with \(\rho_-(t)\to\infty\), such that
\[
\rho_-(d_X(x,x'))\le d_Y(f(x),f(x'))\le\rho_+(d_X(x,x'))
\quad(x,x'\in X).
\]
This does not by definition require the affine-linear control of a quasi-isometric embedding. In a smooth Riemannian setting, a map \(h:X\to Y\) is harmonic when its tension field vanishes, \(\operatorname{tr}\nabla\mathrm dh=0\). For singular metric targets, a suitable energy-based harmonic notion must instead be specified.

A customary strong approximation assertion would produce such an \(h\) with
\[
\sup_{x\in X}d_Y(f(x),h(x))<\infty,
\]
but the source does not state the approximation theorem or its actual error control. It therefore does not assert that every coarse embedding between arbitrary metric spaces admits this approximation.

Hulin and Pansu's question is to find applications of the intended harmonic-map approximation result. This is a programme for deriving new consequences, not a specific proposed estimate or universal existence theorem. The domain and target classes, curvature/completeness hypotheses, approximation notion, and assumptions on \(f\) must be recovered before a candidate application can be formulated precisely. The supplied research summary mentions nonpositively curved targets and emerging rigidity/quasi-isometry applications, but gives no actual consequence or complete theorem to formalize.
\subsection{Short English statement}
Find concrete consequences of the intended theorem approximating coarse embeddings by harmonic maps. The approximation hypotheses and error bound are absent, so no universal approximation statement or particular application is invented.
\subsection{Sources}
\begin{itemize}
\item Local UnsolvedMath registry, permanent ID \(6800013\), exact source code \texttt{catalogue entry}, ``Coarse embeddings --- Question 2.'' Proposers Dominique Hulin and Pierre Pansu, Paris-Sud (Paris-Saclay), 2018.
\item Morgan and Pansu, \emph{A List of Open Problems in Differential Geometry} (2018), Question environment 13, ``Coarse embeddings''; \url{https://www.imo.universite-paris-saclay.fr/~pansu/problems_MTDG.tex}. Source-TeX extraction; recorded access 2026-07-29.
\item Supplied OPEN-TRIAGE research summary describes the Hulin--Pansu programme and partial applications, but supplies no complete approximation-theorem citation or precise application statement.
\end{itemize}
\end{document}
